% Generated draft: requires mathematical and source review.
\documentclass[11pt]{article}
\usepackage[margin=25mm]{geometry}
\usepackage{fontspec}
\setmainfont{Latin Modern Roman}
\usepackage{amsmath,amssymb,mathtools,unicode-math}
\setmathfont{Latin Modern Math}
\usepackage{enumitem,needspace,xurl,hyperref}
\setlength{\parindent}{0pt}
\setlength{\parskip}{5pt}
\setlength{\emergencystretch}{3em}
\newcommand{\N}{\mathbb{N}}
\newcommand{\Z}{\mathbb{Z}}
\newcommand{\Q}{\mathbb{Q}}
\newcommand{\R}{\mathbb{R}}
\newcommand{\C}{\mathbb{C}}
\newcommand{\F}{\mathbb{F}}
\newcommand{\A}{\mathbb{A}}
\newcommand{\PP}{\mathbb{P}}
\newcommand{\E}{\mathbb{E}}
\newcommand{\eps}{\varepsilon}
\newcommand{\dd}{\,\mathrm{d}}
\begin{document}
\section*{UnsolvedMath ID: 701435}
\textbf{Source problem code:} catalogue entry
\subsection{Definitions and mathematical statement}
A group \(G\) is a finitely generated one-relator group if it admits a presentation
\[
G\cong \langle x_1,\ldots,x_n\mid r=1\rangle
=F_n/\langle\!\langle r\rangle\!\rangle,
\]
where \(n\geq 1\), \(F_n\) is the free group on \(x_1,\ldots,x_n\), \(r\in F_n\), and \(\langle\!\langle r\rangle\!\rangle\) denotes the normal closure of \(r\).

A group \(A\) is amenable if there exists a finitely additive probability measure \(\mu:\mathcal P(A)\to[0,1]\) on all subsets of \(A\) that is invariant under left translation. Explicitly, \(\mu(A)=1\), \(\mu(\varnothing)=0\), \(\mu(B\cup C)=\mu(B)+\mu(C)\) for disjoint subsets \(B,C\subseteq A\), and \(\mu(aB)=\mu(B)\) for every \(a\in A\) and \(B\subseteq A\).

A group \(G\) is residually amenable if, for every \(g\in G\setminus\{1_G\}\), there exist an amenable group \(A\) and a homomorphism \(\varphi:G\to A\) such that \(\varphi(g)\neq 1_A\). The group \(A\) and the homomorphism may depend on \(g\).

\textbf{Question (G. N. Arzhantseva, Kourovka Notebook 18.6).} Is every finitely generated one-relator group residually amenable? Equivalently, for every integer \(n\geq 1\), every \(r\in F_n\), and every nonidentity element \(g\in F_n/\langle\!\langle r\rangle\!\rangle\), does there exist a homomorphism to an amenable group whose value at \(g\) is nonidentity?

The supplied registry labels this problem open; that label is supplied curation and has not been independently verified.
\subsection{Short English statement}
Must every nonidentity element of a finitely generated group with one defining relation remain nonidentity under some homomorphism to an amenable group? The target group may depend on the element. The question is attributed to G. N. Arzhantseva.
\subsection{Sources}
\begin{itemize}
\item Local UnsolvedMath registry, supplied complete source record: permanent ID \(701435\), exact source problem code \texttt{catalogue entry}, title \emph{Kourovka Notebook Problem 18.6}. The supplied status is \texttt{open}.
\item \emph{The Kourovka Notebook (Unsolved Problems in Group Theory)}, Problem 18.6, proposed by G. N. Arzhantseva in 2014. The supplied background identifies the 18th issue (2014) and the 21st issue edition (2026), with identifier \texttt{arXiv:1401.0300v48}. The supplied collection description names E. I. Khukhro and V. D. Mazurov as editors. These bibliographic details are reproduced from the record and have not been independently checked.
\end{itemize}
\end{document}
